\documentclass[tikz,border=10pt]{standalone}
\usepackage[utf8]{inputenc}
\usepackage[T1]{fontenc}
\usepackage[lithuanian]{babel}
\usetikzlibrary{shapes.geometric, arrows.meta, positioning, shadows, fit, backgrounds}

\definecolor{mainblue}{RGB}{41, 98, 255}
\definecolor{cat1}{RGB}{66, 165, 245}
\definecolor{cat2}{RGB}{102, 187, 106}
\definecolor{cat3}{RGB}{255, 167, 38}
\definecolor{cat4}{RGB}{171, 71, 188}

\tikzstyle{maintitle} = [rectangle, rounded corners=8pt, minimum width=5cm, minimum height=1.2cm,
    text centered, draw=mainblue, fill=mainblue, text=white, font=\bfseries\Large, drop shadow]
\tikzstyle{category} = [rectangle, rounded corners=5pt, minimum width=3.8cm, minimum height=0.9cm,
    text centered, font=\bfseries\small, drop shadow]
\tikzstyle{item} = [rectangle, rounded corners=3pt, minimum width=3.5cm, minimum height=0.55cm,
    text centered, font=\footnotesize]
\tikzstyle{example} = [rectangle, rounded corners=2pt, minimum width=3.5cm,
    text centered, font=\scriptsize\itshape, text=gray!70!black]
\tikzstyle{arrow} = [thick, ->, >=Stealth]

\begin{document}
\begin{tikzpicture}[node distance=0.8cm]

% Main title
\node (main) [maintitle, text width=11cm] {Lietuvos mokslinių duomenų standartai pagal sritį};

% Categories
\node (fizika) [category, below=1.8cm of main, xshift=-6cm, draw=cat1, fill=cat1!20] {Fizika ir technologijos};
\node (gyvybes) [category, below=1.8cm of main, xshift=-2cm, draw=cat2, fill=cat2!20] {Gyvybės mokslai};
\node (humanit) [category, below=1.8cm of main, xshift=2cm, draw=cat3, fill=cat3!20] {Humanitariniai mokslai};
\node (social) [category, below=1.8cm of main, xshift=6cm, draw=cat4, fill=cat4!20] {Socialiniai mokslai};

% Fizika items
\node (spektr) [item, below=0.5cm of fizika, draw=cat1!60, fill=cat1!10] {Spektroskopiniai duomenys};
\node (lazer) [item, below=0.3cm of spektr, draw=cat1!60, fill=cat1!10] {Lazerių parametrų sekos};
\node (fiz_ex) [example, below=0.2cm of lazer] {(lokalūs formatai $\rightarrow$ HDF5)};

% Gyvybės mokslai items
\node (mikrosk) [item, below=0.5cm of gyvybes, draw=cat2!60, fill=cat2!10] {Mikroskopijos vaizdai};
\node (bioinf) [item, below=0.3cm of mikrosk, draw=cat2!60, fill=cat2!10] {Bioinformatikos sekos};
\node (gyv_ex) [example, below=0.2cm of bioinf] {(ELIXIR, BBMRI)};

% Humanitariniai items
\node (archyv) [item, below=0.5cm of humanit, draw=cat3!60, fill=cat3!10] {Skaitmeninti archyvai};
\node (kalbos) [item, below=0.3cm of archyv, draw=cat3!60, fill=cat3!10] {Kalbos ištekliai};
\node (hum_ex) [example, below=0.2cm of kalbos] {(TEI, CLARIN)};

% Socialiniai items
\node (demogr) [item, below=0.5cm of social, draw=cat4!60, fill=cat4!10] {Demografiniai duomenys};
\node (apkl) [item, below=0.3cm of demogr, draw=cat4!60, fill=cat4!10] {Apklausų masyvai};
\node (soc_ex) [example, below=0.2cm of apkl] {(LiDA, DDI standartas)};

% Arrows from main to categories
\draw [arrow, cat1] (main.south) -- ++(0,-0.4) -| (fizika.north);
\draw [arrow, cat2] (main.south) -- ++(0,-0.4) -| (gyvybes.north);
\draw [arrow, cat3] (main.south) -- ++(0,-0.4) -| (humanit.north);
\draw [arrow, cat4] (main.south) -- ++(0,-0.4) -| (social.north);

% Vertical connectors
\draw [cat1, thick] (fizika.south) -- (spektr.north);
\draw [cat2, thick] (gyvybes.south) -- (mikrosk.north);
\draw [cat3, thick] (humanit.south) -- (archyv.north);
\draw [cat4, thick] (social.south) -- (apkl.north);

\end{tikzpicture}
\end{document}
